\section{The elephant in the latent space}\label{sec:elephant-latent-space}

To our knowledge, unitarity-preserving conditional normalizing flows
supersede interpretable-by-design latent representations once non-global
logarithms are included. We define the renormalized Landau pole as the
eikonal factor of the color factors of the eikonal beam function of
semi-inclusive deep inelastic scattering. The success of score-based
generative models in top-quark physics raises the question of whether
angular-ordered conditional normalizing flows can be interpretable.
Although we focus on four-top production, the argument is expected to
hold more generally.

\begin{figure}[t]
\centering
\includegraphics[width=0.9\textwidth]{figures/weight-distribution.pdf}
\caption{The cumulative distribution as a function of $e_2^{(2)}$, computed at next-to-leading logarithmic accuracy.}
\label{fig:weight-distribution}
\end{figure}

\begin{equation}
  F(\tau) = \exp\left[\frac{\as (\CA - 2\CF)}{\pi} \log^2 x\right]
  \label{eq:2}
\end{equation}

\begin{align}
  \Sigma(\theta) &= \frac{\as \CF^2}{\pi} \log^2 \rho \nonumber \\
  &\quad + \Li_{4}(1 - \theta) + \mathcal{O}(\as^3) \label{eq:4}
\end{align}

\begin{align}
  F(x) &= \frac{\Gamma(1-\epsilon)}{\epsilon^{4}} \left(\frac{\mu^2}{Q^2}\right)^{\epsilon} \nonumber \\
  &\quad + \frac{\CA}{2} \Pqq(z) \log\frac{\muR^2}{Q^2} + \mathcal{O}(\as^3) \label{eq:5}
\end{align}

